%% ============================================================
%% 分片 11：原文第 5 节「Topological Field Theory In Two Dimensions」（原文 glshort.tex 起于第 3264 行）
%% 由 tools/harvmac2tex.py 从 arXiv 源码（hep-th/0604151v3）自动生成。
%% 只译英文散文；公式、编号、\label、\cite 一律不改。
%% ============================================================

\section{Topological Field Theory In Two Dimensions}
\label{sec:tqfttwo}

In section 第~\ref{sec:tqftfour}~节, we analyzed a twisted version of ${\cal
N}=4$ super Yang-Mills theory on a general four-manifold $M$.  The
twisting preserved those supersymmetries that are invariant under
a subgroup $Spin'(4)\subset Spin(4)\times Spin(6)$.  The condition
for a supersymmetry generator to be invariant under $Spin'(4)$ is
% 式 (5.1)
\setcounter{equation}{0}
\begin{equation}\label{eq:tory}
\left(\Gamma_{\mu\nu}+\Gamma_{4+\mu,4+\nu}\right)\epsilon=0,
~~\mu,\nu=0,\dots,3.
\end{equation}
  In addition, all 16 supersymmetry
generators obey a ten-dimensional chirality condition,
% 式 (5.2)
\setcounter{equation}{1}
\begin{equation}\label{eq:nory}
\Gamma_0\Gamma_1\dots\Gamma_9\epsilon=\begin{cases}\epsilon&
~\text{in Lorentz signature}\\-i\epsilon&~\text{in Euclidean
signature.}\\\end{cases}
\end{equation}
 In the twisted theory, four of the six scalar
fields $\phi_0,\dots,\phi_5$ of ${\EUN}=4$ super Yang-Mills theory
are reinterpreted as a one-form
$\phi=\sum_{\mu=0}^3\phi_\mu\,dx^\mu$. Two scalar fields, $\phi_4$
and $\phi_5$, are untwisted, and are rotated by a group
$SO(2)_{\cal R}$ (whose generator we called ${\cal K}$ in section
第~\ref{sec:twisting}~节). Allowing for the fermions, the symmetry group of the
twisted theory is $Spin(2)_{\cal R}$, the double cover of
$SO(2)_{\cal R}$.

The twisted theory is a formal construction in the sense that
twisting violates unitarity and only works in Euclidean
signature.\footnote{To construct a twisted theory in Lorentz
signature, we would have needed a suitable homomorphism
$Spin(1,3)\to Spin(6)$.  Because of the compactness of $Spin(6)$,
a non-trivial homomorphism does not exist. If one replaces
$Spin(6)$ by $Spin(1,5)$ to work around this, the couplings to
fermions will no longer be hermitian and the energy is no longer
bounded from below.} Suppose, however, that we split off the time
direction and take $M=\Bbb{R}\times W$, where $\Bbb{R}$
parametrizes the time, $W$ is a three-manifold, and the metric on
$M$ is  a product. Then time is not involved in the twisting, and
the twisted theory makes sense with Lorentz or Euclidean signature
and is unitary and physically sensible. In this case, as the
holonomy of $M$ reduces to $SO(3)$, unbroken supersymmetries need
only be invariant under a subgroup $Spin'(3)\subset Spin'(4)$;
equivalently, we need only impose \eqref{eq:tory}\ for $\mu,\nu=1,\dots,3$.
There are four unbroken supersymmetry generators, which obey
\eqref{eq:tory}\ for $\mu,\nu=1,\dots 3$. Three scalars are untwisted,
namely $\phi_0$, $\phi_4$ and $\phi_5$, and are rotated by a
$Spin(3)_{\cal R}$ symmetry that extends the $Spin(2)_{\cal R}$
symmetry that is present on any $M$.

Let us specialize this further to the case $M=\Bbb{R}^{1,1}\times
C$ (or $\Bbb{R}^2\times C$), for $C$ a Riemann surface.  The
holonomy group of $M$ reduces to $SO(2)$, and \eqref{eq:tory}\ collapses to
the single condition % [WARN] 公式 \eqn\eqref{eq:ytytag} 无编号或重复
\begin{equation*}
\tag*{(5.3)}
\Gamma_{2367}\epsilon=\epsilon,\label{eq:ytytag}
\end{equation*}

leaving an eight-dimensional space of unbroken supersymmetries.
Four scalars $\phi_0,\phi_1,\phi_4$, and $\phi_5$ are untwisted,
and the symmetry group that rotates the untwisted scalars is now
extended to $Spin(4)_{\cal R}$.

In fact, the structure in two dimensions is precisely that of
$(4,4)$ supersymmetry. If we combine \eqref{eq:ytytag}\ and \eqref{eq:nory}, we find (in
Lorentz signature) % 式 (5.4)
\setcounter{equation}{3}
\begin{equation}\label{eq:byty}
\Gamma_{01}\epsilon
=\Gamma_{4589}\epsilon.
\end{equation}
 Here $\Gamma_{01}$ is the operator that
distinguishes the two spinor representations of the Lorentz group
$SO(1,1)$ of $\Bbb{R}^{1,1}$.
 Similarly, $\Gamma_{4589}$ distinguishes the two spin
representations of $Spin(4)_{\cal R}=SU(2)'\times SU(2)''$. These
two spin representations transform under $SU(2)'\times SU(2)''$ as
$({\bf 2},{\bf 1})$ and $({\bf 1},{\bf 2})$, respectively. Eqn.
\eqref{eq:byty}\ says that the eigenvalues of $\Gamma_{01}$ and
$\Gamma_{4589}$ are equal for all two-dimensional supersymmetries.

So with a suitable choice of orientation of $\Bbb{R}^{1,1}$, the
supersymmetries transform as $({\bf 2},{\bf 1})^-\oplus({\bf
1},{\bf 2})^+$, where the superscript labels the eigenvalue of
$\Gamma_{01}$. But this is the usual structure of $(4,4)$
supersymmetry in two dimensions: one $SU(2)$ group of
$R$-symmetries -- namely $SU(2)'$ -- acts on the supersymmetries
of negative $Spin(1,1)$ chirality, while a second such group --
$SU(2)''$ --  acts on supersymmetries of positive chirality. In
fact, in section 第~\ref{sec:compacthitchin}~节, we have already identified the
effective low energy theory as  a sigma-model in which the target
space is the hyper-Kahler manifold $\MH$.  In general,
sigma-models with hyper-Kahler targets have $(4,4)$ supersymmetry
\cite{ref65}.

In addition to the $Spin(4)_{\cal R}$ global symmetry, the twisted
gauge theory on $\Bbb{R}^{1,1}\times C$ has an additional symmetry
that appears because the holonomy group of $C$, by which we twist,
is abelian.  Even after twisting $\phi_2$ and $\phi_3$, it is
still possible to make a rotation of the $\phi_2-\phi_3$ plane.
This is the symmetry group, called ${\cal U}_1$ in section
第~\ref{sec:mhhyper}~节, that acts on a Hitchin pair $(E,\varphi)$ by rotating
the phase of $\varphi$.
